\section{Conclusion}
\label{sec:conclusion}

\sys{} takes the language model as it is---capable, fallible, and
steerable by what it reads---and builds around it a platform in which
the model holds no authority. Every proposed action meets gates whose
inputs are identity, grants, a person-chosen trust level, an agent's
approved manifest, and the person's approval of the exact call; none is
derived from model output. Untrusted agent code runs as its own user,
fenced by Landlock and seccomp, reaching only declared hosts through a
proxy, with credentials lent per call and never shown to the model.
What appears as data beside the model's words is drawn from stored,
validated results. A stop withdraws the keys rather than asking for
cooperation, and an append-only trail records what was allowed.

None of these mechanisms is new on its own; they are the protections
operating systems have offered programs for half a century. The
contribution is to place them around a planner that is very capable and
not trusted, so that when it is wrong---and it will be---the mistake is
contained to what a person could already do, visible on a card or in a
record, and stoppable. The implementation is open, the tests that hold
its claims can be rerun by anyone, and the limits are stated beside
them.

\section*{Acknowledgements}
AI tools (Claude, by Anthropic) assisted in drafting the text of this
paper and in checking its claims against the source code. The author
reviewed the text and is responsible for its content.

\section*{Availability}
The source code, documentation, and test suites are available at
\url{https://github.com/decentai-io/decentai}; the agent catalog used in
the examples is at \url{https://github.com/decentai-io/decentai-agents}.
